% Reserved D.68 candidate. Unregistered pending independent proof audit.
\subsection{Compactness after fixing the actual abelian cocycle}

\begin{statement}{D.68}{Proved}{Fixed-abelian nilpotent fibres are compact and detect finite failure}
Retain the actual tangential arithmetic action on
\(T_n=F_5/\gamma_{n+1}F_5\). Fix an actual continuous
abelian cocycle
\(\alpha:G_{\mathbb Q}\to T_1=\mathbb Z_5(1)^2\),
with unrestricted ramification. Give continuous cochain
spaces their uniform topology. Then, for every \(n\),
\[
 \mathscr Z_n(\alpha)=
 \{z\in Z^1_{\rm cts}(G_{\mathbb Q},T_n):
                     z\bmod\gamma_2=\alpha\}
                                                        \tag{D.68.1}
\]
is compact. It is the image of a closed subset of a
finite power of \(\mathbb Z_5\) under a continuous
family of actual cocycles whose Hall coordinates
are rational polynomials in the parameters with
coefficients in fixed finite spaces of continuous
cochains. Consequently the section-class fibre
\[
 \mathscr S_n([\alpha])=
 \operatorname{Sec}(T_n)\mathbin{\times}_{\operatorname{Sec}(T_1)}
                                      \{[\alpha]\}
                                                        \tag{D.68.2}
\]
is compact and Hausdorff, as is the analogous full
free pro-five fibre. Every transition map between
these compact fibres is continuous and has closed
image.

If a fixed class \(u\in\mathscr S_m([\alpha])\)
lifts to \(T_n\) for every \(n\ge m\), then these
individual lifts can be chosen compatibly, and
\(u\) lifts to the actual full free pro-five group.
The same assertions hold for the actual nilpotent
quotients \(M_n=M/\gamma_{n+1}M\) of
\(M=F_5/F_5''\), with their limit \(M\).
In particular, individual nonemptiness at every
nilpotent level suffices for nonemptiness of the
corresponding actual limit fibre.

For \(a\) in the generic cusp set of D.67, every
noncuspidal class at a fixed finite nilpotent stage
above \((a,0)\) fails to lift at some higher finite
stage. More precisely, for every neighborhood \(V\)
of the cusp class in \(\mathscr S_m((a,0))\),
there is an \(N\ge m\) such that the image of
\(\mathscr S_N((a,0))\) in
\(\mathscr S_m((a,0))\) is contained in \(V\).
At class three the fibre is the compact torsor
\(\mathbb Z_5^2\), with cusp as origin; for every
\(h\ge1\), one finite stage excludes every
parameter outside \(5^h\mathbb Z_5^2\).
For the metabelian quotients this last assertion
already holds on the generic set of D.66.
These finite stages are not evaluated here; an
explicit first stage and its obstruction class
remain open.
\end{statement}

\begin{proof}
The issue is compactness inside the actual continuous
cochain complex. In particular, the global group
\(H^2\) may be nonseparated, as D.1 and D.62 prove.
We shall use a finite coefficient module inside
that complex and shall never replace the subgroup
of coboundaries by its closure.

\paragraph{Integral coordinates in every finite class.}
Let \(F_{\rm disc}\) be the discrete free group on
\(x,y\). The basic commutator theorem identifies
\(F_{\rm disc}/\gamma_{n+1}F_{\rm disc}\) with the
ordered integral powers of its basic commutators
of weights at most \(n\). Its successive
lower-central factors are free abelian on the
commutators of the corresponding weight; see
\cite[Theorem 5.6 and its corollary]{DE_HallEdmonton}.
In these coordinates multiplication and powers
are integer-valued rational polynomials
\cite[Theorem 6.5]{DE_HallEdmonton}. The same
source immediately after that theorem, case (iii),
identifies the coordinate construction over
\(\mathbb Z_5\) with the pro-five completion.

For clarity, this completion is exactly \(T_n\).
A homomorphism from the discrete free group to
a finite five-group of class at most \(n\)
factors through the displayed discrete quotient,
and extends uniquely from the free pro-five
group. Conversely every finite quotient of
\(T_n\) has class at most \(n\). Thus the two
pro-five completions have exactly the same
finite quotients with the specified generators.
It follows that the Hall coordinate map is a
homeomorphism from a finite power of
\(\mathbb Z_5\) onto \(T_n\), and that
\[
 L_n=\ker(T_n\longrightarrow T_{n-1})
       \simeq\mathbb Z_5(n)^{d_n},\qquad n\ge2,
                                                        \tag{D.68.3}
\]
where \(d_n\) is the number of basic commutators
of weight \(n\). The arithmetic action is scalar
\(\chi_5^n\) on this factor, since it is scalar
\(\chi_5\) on degree one and the graded Lie ring
is generated in degree one. Inner twisting by
any lower section does not change the action
on the central kernel \(L_n\).

Appending zero top Hall coordinates to the
lower coordinates gives a continuous polynomial set-section
\(j_n:T_{n-1}\to T_n\). No assertion of integral
Baker--Campbell--Hausdorff coordinates in class
at least five is needed. The rational denominators
of the Hall polynomials are bounded at each
fixed class, and their values on the actual
coordinate lattice remain integral.

The actual arithmetic action is polynomial in
the Hall coordinates of its input, with coefficients
that are fixed continuous functions of the
arithmetic element. Indeed an automorphism is
determined by its images of \(x,y\). Each basic
commutator maps to the same fixed group word
in these two images, and its arbitrary
\(\mathbb Z_5\)-power is evaluated by the power
polynomials just cited. Substituting these finitely
many coordinates in the multiplication polynomials
proves this assertion at every fixed class.
This describes the genuine action, including
all its higher arithmetic coefficients.

\paragraph{Finite generation of actual central cocycles.}
For every integer \(j\ge2\), every continuous
\(\mathbb Z_5(j)\)-valued one-cocycle kills
inertia at every \(q\ne5\). Its restriction to
inertia is a homomorphism because the coefficient
is unramified. Wild inertia and the prime-to-five
tame factors are killed. If \(b\) is its value
on a five-primary tame generator, the tame
Frobenius relation and the cocycle identity give
\[
                         q b=q^j b.
                                                        \tag{D.68.4}
\]
Since \(q^j-q\ne0\) in \(\mathbb Z_5\), we
have \(b=0\). The coefficient action itself is
unramified there. Thus the actual affine cocycle
factors through \(G_{\mathbb Q,\{5\}}\).

Finite-coefficient cohomology of this last
group is finite
\cite[I, Corollary 4.15]{DE_MilneADT}.
The finite systems of invariant groups are
Mittag--Leffler, so continuous comparison
\cite[Proposition 1.6 and Theorem 2.2]{DE_Jannsen1988}
identifies integral \(H^1\) with the inverse
limit of its finite-coefficient groups. Give it
this inverse-limit topology for the moment; it
is compact in that topology, before any assertion
about the uniform cochain quotient. The coefficient sequence
for multiplication by five injects
\(H^1/5H^1\) into the finite group
\(H^1(G_{\mathbb Q,\{5\}},\mathbb F_5(j))\).
Compact Nakayama then makes integral \(H^1\)
a finitely generated \(\mathbb Z_5\)-module.
This is the finiteness portion of D.1, now
spelled out for every \(j\ge2\); no finite
ramification is imposed on the original cocycle.

Choose finitely many actual continuous cocycles
representing generators of this \(H^1\), and
add the coboundary of a basis of the coefficient
lattice. They generate the module of all actual
one-cocycles. A finite integral linear combination
of these fixed continuous maps is continuous,
and the map from the finite parameter power of
\(\mathbb Z_5\) is continuous for the uniform
cochain topology. Hence
\[
 Z^1_{\rm cts}(G_{\mathbb Q},L_n)
       \text{ is a finitely generated compact }
                          \mathbb Z_5\text{-module}.
                                                        \tag{D.68.5}
\]
Torsion in \(H^1\) causes no problem: its
representative cocycles, together with the
coboundaries, still generate all actual cocycles.
The latter module is a submodule of a torsion-free
cochain module.

\paragraph{The finite-span coboundary lemma.}
Let \(L\) be a finite free \(\mathbb Z_5\)-module
with continuous arithmetic action. Fix finitely
many continuous two-cochains
\(\zeta_i\in C^2(G_{\mathbb Q},L\otimes\mathbb Q_5)\),
for \(1\le i\le r\), and define the actual
coefficient submodule
\[
 B=\left\{v\in\mathbb Z_5^r:
       \sum_i v_i\zeta_i\in
                 d C^1(G_{\mathbb Q},L)\right\}.
                                                        \tag{D.68.6}
\]
It is a \(\mathbb Z_5\)-submodule, even though
the coboundaries in the entire cochain space
need not be closed. Since \(\mathbb Z_5\) is
Noetherian, \(B\) is finite free. It is closed:
the map from \(\mathbb Z_5^e\) given by any
module basis of \(B\) has compact image in
\(\mathbb Z_5^r\).

Choose a basis \(v_1,\ldots,v_e\) of \(B\)
and integral continuous primitives \(b_k\) with
\(d b_k=\sum_i(v_k)_i\zeta_i\). The coordinate
functionals of this basis extend to linear
functionals over \(\mathbb Q_5\) on
\(\mathbb Q_5^r\): choose a left inverse over
\(\mathbb Q_5\) of the matrix with columns
\(v_k\). Consequently, if \(v=v(t)\) is a
polynomial coefficient vector taking values
in \(B\), the primitive
\[
                b(t)=\sum_k (R v(t))_k b_k
                                                        \tag{D.68.7}
\]
is integral and continuous, and is again a
rational-polynomial family in a fixed finite
space of continuous one-cochains. The coordinates
\((R v(t))_k\) are integral on the allowed
parameter set. Their denominators as polynomial
functions are fixed and bounded. This is an
actual choice of primitives, not a lifting
through a separated cohomology group.

\paragraph{Induction with all exact cocycles.}
We prove the stronger finite-parameter assertion
in the statement by induction on \(n\). A
polynomial family means a family on a closed
set \(X\subset\mathbb Z_5^r\), for a finite
\(r\), such that every Hall coordinate has
an expression
\[
                \sum_i P_i(t)\varphi_i(g),
 \quad P_i\in\mathbb Q_5[t_1,\ldots,t_r],\quad
       \varphi_i\in C(G_{\mathbb Q},\mathbb Q_5).
                                                        \tag{D.68.8}
\]
Only finitely many polynomials and fixed functions
occur. The coordinate values are integral on
\(X\). Since both parameter space and Galois
group are compact, joint continuity gives
continuity into the uniform cochain space.
At \(n=1\), the single family member is the
fixed actual abelian cocycle \(\alpha\).

Suppose that a polynomial family \(f_t\),
parametrized by a closed \(X\), covers all
of \(\mathscr Z_{n-1}(\alpha)\). Apply \(j_n\)
to obtain actual set-theoretic lifts \(\widehat f_t\).
Their multiplication defects are integral
continuous central two-cocycles in \(L_n\):
\[
 \eta_t(g,h)=\widehat f_t(g)\,
        \sigma_0(g)(\widehat f_t(h))\,
                   \widehat f_t(gh)^{-1}.
                                                        \tag{D.68.9}
\]
The coordinate facts proved above show that
this is a polynomial family in a fixed finite
space of two-cochains. To make the finite-span
claim explicit, multiplication and inversion
use finitely many Hall polynomials; arithmetic
substitution uses the finitely many continuous
coordinates of \(\sigma_0(g)(x)\) and
\(\sigma_0(g)(y)\); and substitution of
\(f_t(g),f_t(h),f_t(gh)\) uses finitely many
functions in (D.68.8). Expanding their finite
products gives finitely many monomials in
\(t\), whose coefficients are fixed continuous
functions of \((g,h)\). This proves the
finite-span assertion uniformly on all of \(X\),
including after every previous choice of
primitives.

Absorb bounded polynomial denominators into
the coefficient cochains. We may therefore
write
\[
                 \eta_t=\sum_i v_i(t)\zeta_i,
       \qquad v(t)\in\mathbb Z_5[t_1,\ldots,t_r]^k,
                                                        \tag{D.68.10}
\]
where the fixed \(\zeta_i\) may be rational
rather than integral cochains. Form \(B\) as
in (D.68.6), with the actual integral coefficient
lattice \(L_n\). Exactly those parameters in
\[
                         X'=X\cap v^{-1}(B)
                                                        \tag{D.68.11}
\]
admit an integral continuous cocycle lift.
This set is closed and compact. Choose the
primitive \(b(t)\) of (D.68.7); multiplying
\(\widehat f_t\) by the central cochain
\(-b(t)\) cancels its multiplication defect
and produces a genuine lifted cocycle.

Choose finite generators \(z_1,\ldots,z_e\)
of the actual central cocycle module (D.68.5).
Multiplying the chosen lift by
\(\sum_j\lambda_j z_j\), with
\(\lambda\in\mathbb Z_5^e\), gives a
polynomial family on the compact closed set
\(X'\times\mathbb Z_5^e\). It covers every
exact lift: the difference of any two lifts
with identical projection is an actual central
one-cocycle. By the inductive surjectivity onto
all lower exact cocycles, the new family covers
all of \(\mathscr Z_n(\alpha)\).
The chosen primitives are rational-polynomial
in the old parameters, and the new cocycles
are linear in the new parameters. Thus the
finite-span property (D.68.8) is preserved,
not merely continuity or a statement about
cohomology classes. This finishes the induction
and proves compactness in (D.68.1).

\paragraph{Conjugacy, limits, and closed images.}
The uniform cocycle space is Hausdorff, and
conjugation by the compact group \(T_n\) has
closed orbits and a Hausdorff quotient. Every
class with abelian class \([\alpha]\) admits
an exact representative with abelian cocycle
\(\alpha\): lift the required abelian
conjugator from \(T_1\) to \(T_n\).
A conjugator between two such exact
representatives has abelian image fixed by
\(G_{\mathbb Q}\), hence zero. Thus
\(\mathscr S_n([\alpha])\) is also the
quotient of (D.68.1) by the compact group
\(\gamma_2T_n\). It is compact Hausdorff.
Projection of cochains is continuous, so the
transition maps have closed image.

Now fix an actual representative \(z_m\) of
\(u\). At each \(n\ge m\), its exact lift
space is a closed subset of
\(\mathscr Z_n(\alpha)\), hence compact.
If \(u\) has a class lift at that level,
lifting a conjugator from \(T_m\) makes
this exact lift space nonempty. These compact
spaces form an inverse system. Every finite
compatibility condition can be met by a
member at the highest of the finitely many
levels. Compactness therefore gives a
compatible sequence of actual cocycles.
Since \(F_5=\varprojlim T_n\), their inverse
limit is an actual continuous \(F_5\)-cocycle.
The uniform cochain topology into \(F_5\) is
exactly the inverse-limit topology on these
exact fibres: every finite quotient of \(F_5\)
is nilpotent and therefore factors through
some \(T_n\), so every uniform neighborhood
is detected at a finite stage.
The same compact construction with \(m=1\)
proves compactness of the full exact fibre;
quotienting by the compact commutator subgroup
proves compactness of the full class fibre.
As in D.67, nested nonempty closed conjugator
sets give injectivity on the inverse limit
of section classes.

The metabelian assertion has exactly the same
hypotheses. To check them, the discrete free
metabelian Alexander module is
\(\mathbb Z[X^{\pm1},Y^{\pm1}]C\), and
its lower-central filtration is given by
powers of \((X-1,Y-1)\). Its nilpotent
quotients are therefore finitely generated
torsion-free nilpotent groups, with central
grade \(n\) free of rank \(n-1\).
Their pro-five completions are the actual
\(M_n\), by the same finite-quotient universal
property; equivalently use the completed
Alexander description of D.35. The central
grade has arithmetic action \(\chi_5^n\).
Hall's polynomial theorem applies to these
nilpotent groups too. The preceding proof
therefore supplies compact exact fibres and
compatible representatives. Here the limit is
indeed \(M\): the intersection of all augmentation
ideal powers in \(\mathbb Z_5[[u,v]]\) is zero,
and compactness then makes the map to the
inverse limit of the quotients surjective.
This also removes a possible compatibility
ambiguity in the individual-stage version
of the metabelian branch criterion of D.41.

\paragraph{The generic finite-stage obstruction.}
For the parameters of D.67, the intersection
of the descending compact images
\[
 \operatorname{im}\bigl(\mathscr S_N((a,0))
                      \longrightarrow\mathscr S_m((a,0))\bigr),
                         \qquad N\ge m,
                                                        \tag{D.68.12}
\]
is exactly the cusp class. Indeed a class
in their intersection lifts at every stage,
hence to \(F_5\), and D.67 identifies that
full fibre. If an open neighborhood \(V\)
of the cusp failed to contain every image
for large \(N\), the compact sets obtained
by intersecting (D.68.12) with the complement
of \(V\) would all be nonempty and nested.
Their intersection would contradict the
singleton intersection just proved.

For \(m=3\), the central kernel over \(T_2\)
is \(\mathbb Z_5(3)^2\). The exact fibre
above its cusp lift is a torsor under
\(H^1(G_{\mathbb Q},\mathbb Z_5(3))^2
=\mathbb Z_5^2\), by D.1 and the central
lifting calculation. There is no stabilizer
quotient: a fixed element of \(T_2\) for a
section action has a first nonzero degree
with positive cyclotomic weight, which is
impossible. Also the class-two fibre above
\((a,0)\) is a singleton, since
\(H^1(G_{\mathbb Q},\mathbb Z_5(2))=0\).
Explicitly, twist D.1's residual unit module
\(\mathbf1\oplus\omega\oplus\omega^2\) by
\(\omega\). The resulting characters
\(\omega,\omega^2,\omega^3\) have no invariant
part, so \(H^1(G_{\mathbb Q,\{5\}},\mathbb F_5(2))=0\).
Compact Nakayama gives the integral vanishing,
and (D.68.4) gives its unrestricted form.
The continuous torsor parametrization is a
bijection from a compact space to the
Hausdorff fibre, hence a homeomorphism.
Taking \(V=5^h\mathbb Z_5^2\) proves the
stated uniform finite failure. The same
argument with the singleton metabelian
fibre of D.66 proves its metabelian version.
No estimate for \(N\), or value of the
first obstruction, follows from this
compactness proof.
\end{proof}
